The shared model receives only normalized operating settings, observed sensors,
availability and the observed-history mask. Engine IDs, subset labels, future
readings, true RUL and simulated altered-entry masks are excluded from inference.
Related-channel fault histories are prefix-consistent: appending future rows does
not alter an earlier corruption realization or select onset by eventual failure.

The four original training files define the engine-disjoint fit/selection/
calibration/internal-test groups in Table~\ref{tab:pooled_partitions}. Scaler,
raw-unit fault scales and repair support use fitting engines only. Fitting-engine
OOF features were used to train the source corrector and exclude their held fold
from temporary fitting and inner selection; in format 10 that corrector is
disabled with a zero bound and contributes no inference feature or adjustment.
Predictor refinement uses fitting engines; every clean/source/NASA, mixed/stuck
and tail-selection check uses held selection engines. Calibration engines enter
only after that choice is fixed.

New artifacts use fresh directories and retain their source SHA-256. Resuming an
interrupted refinement requires matching source, recipe and implementation.
Current detector audits bind the actual deployed thresholds and detector state;
the earlier detector operating point remains historical metadata. Recovery
proofs bind the full declared context matrix, selection IDs and sensor-repair
state. Calibration binds the final model/configuration and its separate scale
scope. Strict reload rejects stale or mismatched records.

Validation-only development runs block official test/RUL and saved test-result
reads. These checks document within-run separation. Earlier internal and
official test cases have already informed development; their later evaluation
remains exploratory and is not an untouched independent generalization test.
